% Gerado a partir de MONOGRAFIA_PASE.md. O texto foi transposto sem alterações.

\chapter{Software na Era da IA Agêntica: Impactos e Lacunas}\label{cap:2}

A Arquitetura PASE não nasce de uma preferência técnica, mas de uma leitura do momento que o ecossistema de software atravessa. Este capítulo apresenta essa leitura. Ele percorre as cinco dimensões nas quais a IA agêntica produz efeitos mais visíveis, a saber, (i) indústria e aplicações, (ii) desenvolvimento de software, (iii) usuário e experiência, (iv) segurança, ética e governança e (v) mercado de trabalho e competências. Em cada uma, o objetivo não é catalogar tendências, mas identificar a lacuna que a separa de uma adoção confiável. A Seção~\ref{sec:2.6} reúne essas lacunas e as converte nos requisitos que orientam a proposta do Capítulo~\ref{cap:4}.

\section{Indústria e Aplicações}\label{sec:2.1}

A história do software empresarial é uma sucessão de modelos que, a cada ciclo, redefiniram não apenas a tecnologia, mas a natureza do que se vende. \citeonline{vista2026outlook} descrevem essa trajetória em três estágios: o software como produto, vendido por licença e instalado localmente; o software como serviço, entregue pela nuvem e cobrado por usuário; e o software como trabalhador, no qual o valor passa a ser medido pelas tarefas executadas. A mudança em curso é de outra natureza. Segundo os autores, a IA agêntica transforma o software empresarial de ferramenta que apoia o trabalho em sistema que executa o trabalho em nome das pessoas.

A escala econômica dessa transição ajuda a explicar a urgência com que as organizações a perseguem. O investimento corporativo em IA triplicou em um único ano e chegou a US\$ 37 bilhões em 2025, e a camada de aplicação concentrou mais da metade desses recursos (MENLO VENTURES, 2025 \textit{apud} SAROYA; MACNEILL, 2026)\nocite{vista2026outlook}. Projeções da Bain \& Company estimam que entre US\$ 5 e 7 trilhões em valor acumulado devem fluir para a camada de software e aplicações até 2030. Trata-se de uma expansão superior a quatro vezes em relação aos níveis atuais (BAIN \& COMPANY, 2024 \textit{apud} SAROYA; MACNEILL, 2026)\nocite{vista2026outlook}, conforme a Figura~\ref{fig:1}.

\begin{figure}[htb]
\centering
\caption{Criação de valor acumulado projetada impulsionada por IA}
\label{fig:1}
\fbox{\parbox[c][4cm][c]{0.85\textwidth}{\centering\small [Imagem: gráfico da Bain \& Company reproduzido em \citeonline{vista2026outlook}]}}
\fonte{\citeonline{vista2026outlook}, a partir de projeções da Bain \& Company (2024).}
\end{figure}

A adoção individual acompanha esse movimento. O relatório do programa DORA (\textit{DevOps Research and Assessment}) de 2025, baseado em 4.867 profissionais de tecnologia de mais de cem países, mostrou que 90\% dos respondentes já usam IA no trabalho, alta de 14,1\% sobre o ano anterior \cite{dora2025}. Os casos empresariais confirmam que a capacidade técnica já não é hipotética. Na Rakuten, um agente de programação implementou uma alteração em uma biblioteca de 12,5 milhões de linhas em sete horas de trabalho autônomo, com 99,9\% de precisão numérica em relação ao método de referência. Na Zapier, a adoção de IA chegou a 89\% da organização, com mais de oitocentos agentes em uso interno \cite{anthropic2026trends}.

Os casos mais maduros, porém, têm um traço em comum que vai além da autonomia: a estrutura. A plataforma Uma, da Duck Creek Technologies, coordena agentes especializados que executam as etapas da subscrição de seguros dentro de um ambiente estruturado e controlado, por onde passam mais de US\$ 60 bilhões em prêmios por ano \cite{vista2026outlook}. Na Fountain, a orquestração hierárquica de agentes acelerou a triagem de candidatos em 50\% e a integração de novos funcionários em 40\% \cite{anthropic2026trends}. \citeonline{vista2026outlook} sintetizam o padrão: a IA gera recomendações probabilísticas, mas o valor desproporcional surge quando as ações finais são executadas de forma determinística, dentro de regras, aprovações e trilhas de auditoria definidas.

O contraste entre a capacidade demonstrada e o resultado capturado define a primeira lacuna. Apenas cerca de 15\% das organizações usam a IA para redesenhar de fato como o trabalho é feito, e ganhos de produtividade na escala da tarefa não têm se convertido em impacto na escala da empresa \cite{wef2026organizational}. Na pesquisa da SPARK (2026)\nocite{spark2026cio}, 54\% dos líderes não conseguem atribuir os custos do processamento autônomo a valor de negócio, e apenas 14\% dos diretores financeiros relatam impacto claro e mensurável. O mesmo levantamento aponta a fragmentação entre funções como a principal barreira para 58\% das organizações, e nomeia o risco que dela decorre: a proliferação descontrolada de agentes em silos isolados, o \textit{Agent Sprawl}.

O mercado ainda não estabilizou sequer a definição do que é um agente. A Information Services Group observa a prática de \textit{agent-washing}, em que soluções convencionais são rebatizadas como agênticas, e registra que apenas 25\% das soluções atuais permitem que agentes operem de forma independente, enquanto 45\% os posicionam como conselheiros de decisores humanos \cite{isg2025agentic}. Os autores concluem que o papel da supervisão humana permanece vagamente definido e que os provedores mais avançados usam a orquestração como camada estruturada de controle, que impõe governança em vez de ser ela mesma autônoma.

A dimensão industrial revela, portanto, uma convergência inédita de capacidade técnica, escala de investimento e pressão competitiva, mas também a ausência de uma arquitetura que transforme essa capacidade em valor governável. O que falta não é mais um agente, e sim um plano de controle que coordene agentes, preserve a decisão humana nos pontos críticos, torne o custo observável e deixe rastro de cada ação.

\section{Desenvolvimento de Software}\label{sec:2.2}

A engenharia de software sempre buscou organizar a complexidade de um processo essencialmente humano. Dos modelos sequenciais às metodologias ágeis, cada paradigma reorganizou as etapas do ciclo de vida, mas manteve a pessoa como agente central de decisão, criação e responsabilidade. Os modelos de linguagem de grande escala (\textit{Large Language Models}, LLMs) rompem esse padrão porque passam a produzir artefatos em praticamente todas as fases do ciclo, dos requisitos à manutenção \cite{burte2025sdlc,zhang2026survey}.

A pesquisa acadêmica reflete essa expansão, mas de forma muito desigual. A revisão sistemática de \citeonline{hou2024llm4se}, com 395 artigos publicados entre 2017 e 2024, mostra que o desenvolvimento de código concentra 56,65\% dos estudos, seguido da manutenção, com 22,71\%, e da garantia de qualidade, com 15,14\%. A engenharia de requisitos responde por apenas 3,9\%, o projeto de software por 0,92\% e a gestão de software por 0,69\%. O levantamento de \citeonline{zhang2026survey}, com 926 estudos classificados em 112 tarefas, confirma o desequilíbrio e conclui que a aplicação de LLMs a requisitos, projeto e gestão continua relativamente pouco explorada.

Esse desequilíbrio tem consequências práticas. A fase em que a IA mais avança, a escrita de código, depende de decisões tomadas antes dela, sobre o que construir e por quê, justamente onde a IA menos foi estudada. \citeonline{hou2024llm4se} observam ainda que, quando a intenção é ambígua, os modelos tendem a produzir código sintaticamente correto e funcionalmente errado, e que apenas seis dos 395 estudos analisados usaram dados industriais. A automação acelera a etapa mais visível do ciclo e deixa a qualidade das decisões que a antecedem por conta de processos que não acompanharam essa aceleração.

A qualidade estrutural do código gerado é a segunda frente de preocupação. \citeonline{kravchuk2025modular} mostram, em experimentos controlados, que um assistente de programação violou o encapsulamento ao acessar membros privados de um módulo em vez de usar sua interface pública, criando dependências ocultas entre módulos. Os autores descrevem o fenômeno como mimetismo arquitetural sem aderência aos princípios subjacentes: o código parece modular, mas não é. Soluções que funcionam na superfície tornam a base de código frágil, e o custo dessa fragilidade aparece apenas mais tarde, na manutenção.

A dificuldade se intensifica quando a tarefa é real. No SWE-bench, que reúne 2.294 problemas extraídos de repositórios públicos em Python, o melhor modelo avaliado resolveu apenas 1,96\% dos casos, em grande parte porque os modelos têm dificuldade em localizar o trecho de código que precisa ser alterado em meio a milhares de arquivos \cite{jimenez2024swebench}. Ainda que os modelos tenham evoluído desde então, o achado permanece instrutivo: resolver um problema real exige saber onde mexer, e esse é um problema de planejamento antes de ser um problema de geração.

Os dados de uso confirmam que a percepção de ganho nem sempre corresponde à medição. Mais de 80\% dos respondentes do relatório DORA acreditam que a IA aumentou sua produtividade, mas 30\% relatam pouca ou nenhuma confiança no código gerado \cite{dora2025}. O mesmo relatório cita um estudo em que desenvolvedores 19\% mais lentos com o uso de ferramentas de IA acreditavam ter ficado 20\% mais eficientes. No plano organizacional, a adoção de IA passou a melhorar o volume de entregas em 2025, mas continua associada a maior instabilidade na entrega de software, o que sugere que os sistemas subjacentes ainda não evoluíram para gerenciar com segurança o desenvolvimento acelerado por IA.

A tese central do relatório DORA é que a IA atua como amplificadora: ela magnifica as forças das organizações de alto desempenho e as disfunções das demais \cite{dora2025}. Entre quinze capacidades testadas, sete mostraram efeito moderador sobre os ganhos, entre elas um ecossistema de dados saudável, práticas sólidas de controle de versão, trabalho em pequenos lotes e foco no usuário. Na ausência dessas capacidades, os ganhos podem ser completamente neutralizados. A fricção, observa o relatório, não desaparece; ela se desloca do trabalho manual para a decisão e a verificação.

A orquestração de múltiplos agentes, apontada como a próxima fronteira do desenvolvimento \cite{anthropic2026trends}, não elimina esse problema e pode agravá-lo. \citeonline{he2025lma} relatam que um sistema multiagente não conseguiu desenvolver de forma autônoma um jogo simples de Tetris e argumentam que a decomposição de requisitos e a gestão de interdependências se tornam mais difíceis à medida que a complexidade cresce. \citeonline{kim2026scaling}, em uma avaliação controlada de arquiteturas de coordenação, encontraram efeitos que variam de ganhos de 80,8\% a perdas de 70\% conforme a tarefa. Em problemas de engenharia de software do SWE-bench Verified, ainda que em uma amostra reduzida, todas as variantes multiagente ficaram abaixo do agente único. Os autores mostram também que agentes que trabalham de forma independente, sem verificação, amplificam erros em até 17,2 vezes, enquanto uma coordenação centralizada, que valida os resultados antes de agregá-los, reduz essa amplificação para 4,4 vezes.

O desenvolvimento de software vive, portanto, uma transformação genuína, mas incompleta. A velocidade de geração cresceu, enquanto a qualidade das decisões anteriores ao código, a integridade estrutural do que é gerado e a verificação do que é entregue não acompanharam o mesmo ritmo. A lacuna não está na capacidade de escrever código, mas na ausência de um processo que submeta essa capacidade a planejamento explícito, verificação determinística e decisão humana informada.

\section{Usuário e Experiência}\label{sec:2.3}

A história da interação humano-computador pode ser lida como uma sequência de interfaces cada vez mais próximas da intenção humana. \citeonline{mukti2025adaptive}, em uma análise bibliométrica de vinte mil documentos publicados entre 1980 e 2024, identificam três fases da personalização. A primeira é uma fase técnica, baseada em regras e filtragem colaborativa, e a segunda, uma fase de escalabilidade, marcada pelo aprendizado profundo. A terceira é uma fase ética, iniciada em 2020, em que a IA generativa permite adaptação em tempo real e a pesquisa passa a priorizar transparência, privacidade e controle do usuário.

A fronteira atual é a das interfaces que deixam de ser artefatos fixos. \citeonline{arora2025selfevolving} descrevem sistemas que se autoevoluem, expandindo o próprio conjunto de funcionalidades em tempo de execução a partir da intenção imediata do usuário. Em um estudo com 72 participantes usando um protótipo bancário, a versão adaptativa obteve 84,38 pontos na Escala de Usabilidade do Sistema, contra 53,96 da versão determinística, uma diferença grande e estatisticamente significativa. Os autores atribuem o ganho à redução do esforço de navegação: em vez de o usuário aprender o modelo mental do sistema, o sistema passa a aprender o do usuário.

O mesmo trabalho, contudo, enumera os riscos dessa abordagem. Uma interface que muda constantemente pode desorientar, porque as pessoas dependem de memória espacial e de leiautes consistentes. Mudanças frequentes aumentam a carga cognitiva, modelos generativos tendem a produzir elementos visualmente inconsistentes e os métodos tradicionais de teste não cobrem o espaço de configurações possíveis \cite{arora2025selfevolving}. Como resposta, os autores recomendam preservar âncoras visuais imutáveis, como a navegação global, e manter os dados transacionais intocados, independentemente de como a apresentação evolui.

A adaptação também precisa ser percebida para ter valor. \citeonline{jamali2025adaptive}, ao comparar três plataformas educacionais com 23 participantes, constataram que as funcionalidades adaptativas foram percebidas como sutis e pouco impactantes, e que a interatividade central das plataformas pesou mais na experiência do que a IA. Os autores identificam um compromisso entre eficiência e engajamento, mediado pela visibilidade das adaptações, e recomendam torná-las explícitas, anunciando ao usuário quando e por que o sistema mudou. A personalização invisível, concluem, não gera valor percebido.

A confiança é a variável que conecta essas evidências. \citeonline{arora2025selfevolving} afirmam que usuários desconfiam de interfaces que mudam sem explicação e que é crítico informar por que cada mudança acontece. \citeonline{bieniek2024generative} apontam a natureza de caixa-preta dos modelos como fonte de desconfiança, especialmente em domínios sensíveis como finanças. Na revisão de \citeonline{tan2025alp}, com 127 estudos sobre plataformas adaptativas, as visualizações que tornavam transparentes as recomendações aumentaram a confiança dos usuários. \citeonline{mukti2025adaptive} propõem mecanismos como um botão que explique cada recomendação.

A personalização tem ainda um custo de privacidade e de interpretação. A precisão da adaptação depende da coleta de dados comportamentais que podem causar desconforto, e a literatura sobre o tema dá pouca atenção à privacidade \cite{tan2025alp}. \citeonline{arora2025selfevolving} defendem que o usuário seja informado sobre o uso de seus dados e possa recusar ou controlar o nível de personalização. \citeonline{tan2025alp} alertam também que um sinal comportamental isolado é ambíguo: uma fixação longa do olhar pode indicar dificuldade, e não engajamento, o que exige combinar várias características antes de adaptar.

A dimensão do usuário mostra que a personalização em tempo real já entrega ganhos mensuráveis de usabilidade, mas expõe uma tensão entre adaptação e previsibilidade, entre conveniência e transparência. A lacuna não é técnica, mas de governança da interface: falta uma forma de personalizar que seja delimitada, explícita para o usuário, reversível ao comportamento padrão e incapaz de introduzir no produto código que não foi revisado.

\section{Segurança, Ética e Governança}\label{sec:2.4}

A segurança de software evoluiu por camadas sucessivas de proteção contra acesso indevido a dados e sistemas. Na IA agêntica, surge uma descontinuidade qualitativa: o adversário não precisa roubar informações, mas pode manipular o próprio processo de decisão do sistema. O relatório anual da CrowdStrike (2026)\nocite{crowdstrike2026threat} abre com a constatação de que o mundo opera na era agêntica, e de que o adversário também. Em 2025, os ataques conduzidos por adversários habilitados por IA cresceram 89\% em relação ao ano anterior, e a IA acelerou o \textit{phishing} e o reconhecimento automatizado, encurtando o tempo entre o acesso inicial e o impacto.

O mesmo relatório registra que os próprios sistemas de IA passaram a integrar a superfície de ataque das organizações, incluindo modelos, dados de treinamento e agentes \cite{crowdstrike2026threat}. Em um caso relatado, adversários embutiram instruções ocultas em um \textit{e-mail} de \textit{phishing} para confundir a triagem automatizada por IA, embora a técnica ainda não tenha demonstrado eficácia consistente em escala. A vulnerabilidade dos agentes é estrutural: um estudo citado por \citeonline{chhabra2026security} encontrou 94,4\% dos agentes avaliados vulneráveis à injeção de instruções (\textit{prompt injection}), 83,3\% a \textit{backdoors} baseados em recuperação e 100\% à exploração da confiança entre agentes (LUPINACCI \textit{et al.}, 2025 \textit{apud} CHHABRA \textit{et al.}, 2026)\nocite{chhabra2026security}.

O código gerado por IA é outro vetor de risco. Modelos treinados com dados de baixa qualidade podem gerar código inseguro, e ataques que contaminam os dados de treinamento podem induzir vulnerabilidades sutis no código produzido \cite{hou2024llm4se}. Quando agentes executam ações, e não apenas recomendam, a defesa recomendada pela literatura combina isolamento da execução, menor privilégio e modelos de proteção que examinam as entradas antes que alcancem o sistema. \citeonline{chhabra2026security} ressaltam, porém, que o isolamento é mais eficaz quando combinado a defesas complementares, e que confirmações humanas também podem ser alvo de engenharia social ou de aprovações desatentas.

A governança de sistemas autônomos é o ponto mais frágil desse quadro. A Information Services Group resume a mudança em uma frase: o que é fundamentalmente novo é a necessidade de governar a ação, já que cada agente capaz de agir sozinho introduz um risco operacional que os modelos atuais de governança ainda não sabem administrar \cite{isg2025agentic}. A rastreabilidade também se torna mais difícil. Uma vez processada por um modelo, a informação é convertida em representações latentes distribuídas, que eliminam os identificadores capazes de ligar uma ação à sua origem e tornam inviável o rastreamento convencional de fluxo de informação \cite{chhabra2026security}.

A dimensão ética acrescenta o problema da responsabilização. Abou Ali e Dornaika (2025)\nocite{abouali2025agentic} observam que a pergunta sobre quem responde por um dano é especialmente aguda em sistemas agênticos baseados em redes neurais, cuja natureza difusa e estocástica desafia marcos jurídicos construídos sobre causalidade direta e intenção. Os autores identificam um padrão por paradigma: sistemas simbólicos predominam em domínios críticos de segurança, como a saúde, enquanto sistemas neurais prevalecem em ambientes adaptativos e ricos em dados, como as finanças. A conclusão é que o caminho mais produtivo é híbrido, com módulos simbólicos garantindo restrições e conformidade sobre decisões baseadas em dados.

A segurança, a ética e a governança de sistemas agênticos não podem ser tratadas como problemas exclusivamente técnicos, nem resolvidas por uma única camada de proteção. A lacuna é a ausência de arquiteturas que tornem a ação dos agentes contida, verificável e atribuível: execução isolada, contratos que limitem o que pode ser produzido, decisões humanas nos pontos em que a ação gera consequência e registros que permitam reconstruir quem decidiu o quê.

\section{Mercado de Trabalho e Competências}\label{sec:2.5}

A pergunta mais frequente sobre IA e trabalho, a da substituição, recebe da literatura analisada uma resposta mais nuançada do que o debate público sugere. Os desenvolvedores usam IA em cerca de 60\% do seu trabalho, mas relatam conseguir delegar por completo apenas entre 0\% e 20\% das tarefas, e cerca de 27\% do trabalho assistido por IA corresponde a tarefas que não teriam sido feitas de outra forma \cite{anthropic2026trends}. O papel do engenheiro se expande em vez de desaparecer: ser engenheiro passa a significar orquestrar agentes que escrevem código, avaliar sua produção, dar direção estratégica e garantir que o sistema resolva o problema certo.

O deslocamento atinge sobretudo a estrutura dos papéis. Na pesquisa da SPARK (2026)\nocite{spark2026cio}, cerca de metade das organizações aponta a capacitação de funcionários para gerenciar, auditar e colaborar com agentes como sua principal adaptação da força de trabalho, e mais de 46\% estão abandonando escopos fixos de cargos em favor de uma evolução contínua de papéis. O Fórum Econômico Mundial registra o surgimento de funções como dono de produto de IA, arquiteto de fluxos de trabalho e curador de modelos, e descreve a média gerência migrando da supervisão de tarefas para a orquestração de fluxos híbridos entre pessoas e IA \cite{wef2026organizational}.

Essas mudanças vêm acompanhadas de tensão. A SPARK (2026)\nocite{spark2026cio} identifica a ansiedade de deslocamento como tema persistente e descreve as respostas das organizações mais bem-sucedidas. Elas reinvestem parte das economias da automação na capacitação das pessoas, dão aos funcionários visibilidade sobre os agentes implantados e reservam cerca de 20\% do tempo para explorações que a IA não consegue simular. O relatório do Fórum acrescenta que o planejamento da força de trabalho deve considerar a capacidade combinada de pessoas e agentes, incluindo a nova carga de revisão e de tratamento de exceções que os próprios agentes criam \cite{wef2026organizational}.

O risco mais sutil está na formação de competências. Em um experimento controlado sobre o aprendizado de uma nova biblioteca de programação, o grupo que usou assistência de IA obteve, em média, 17\% menos pontos na avaliação de conhecimento do que o grupo de controle \cite{shen2026skill}. A diferença dependeu do modo de uso: participantes que delegaram a geração de código ao assistente ficaram entre 24\% e 39\% na avaliação, enquanto os que usaram a IA para perguntas conceituais ou pediram explicações junto ao código ficaram entre 65\% e 86\%. Os autores alertam que profissionais podem não desenvolver as habilidades necessárias para validar e depurar o código escrito pela IA.

O relatório DORA chega a uma conclusão convergente por outro caminho. Os padrões de uso mais comuns da IA estão entregando ganhos expressivos de produtividade e, ao mesmo tempo, bloqueando o desenvolvimento de habilidades da maioria dos desenvolvedores, porque a automação reduz as oportunidades de aprendizado que antes cabiam aos iniciantes \cite{dora2025}. Há aqui um paradoxo de longo prazo: a supervisão humana, que todas as dimensões anteriores apontam como indispensável, depende de pessoas capazes de julgar o que a IA produz, e essa capacidade é exatamente o que o uso desatento da IA corrói.

A dimensão do trabalho mostra que o futuro mais provável não é o da substituição em massa, mas o da redefinição de papéis em torno de orquestração, verificação e julgamento. A lacuna é a de estruturas que façam da supervisão humana um trabalho qualificado e informado, e não uma formalidade de aprovação, e que preservem, no próprio fluxo de trabalho, o engajamento cognitivo de quem decide.

\section{Síntese: Lacunas e Requisitos para uma Arquitetura}\label{sec:2.6}

Tomadas em conjunto, as cinco dimensões revelam um padrão consistente. Em todas elas, a capacidade técnica já demonstrada contrasta com a ausência de estruturas que a tornem confiável. A indústria tem agentes, mas não tem um plano de controle. O desenvolvimento tem geração de código, mas não tem verificação proporcional à velocidade. A experiência do usuário tem adaptação, mas não tem governança da interface. A segurança tem ameaças novas, mas não tem contenção arquitetural da ação. O trabalho tem novos papéis, mas não tem a estrutura que faça da supervisão um ato qualificado.

Essas lacunas não são independentes. A proliferação de agentes sem coordenação amplia a superfície de ataque. A verificação insuficiente do código transfere ao revisor humano uma carga que ele nem sempre tem competência ou contexto para sustentar. A personalização sem transparência corrói a confiança que a adoção exige. Por isso, tratá-las isoladamente, com uma ferramenta para cada sintoma, tende a reproduzir a mesma fragmentação que as originou.

Este trabalho interpreta esse conjunto como a evidência de uma lacuna central: a falta de uma arquitetura que equilibre a autonomia técnica dos sistemas de IA com supervisão humana efetiva, transparência e rastreabilidade, ao longo de todo o ciclo de vida do software. Dessa interpretação derivam cinco requisitos para uma arquitetura que pretenda preencher essa lacuna, apresentados no Quadro~\ref{quadro:1} e retomados no Capítulo~\ref{cap:4}.

\begin{quadro}[htb]
\centering
\caption{Lacunas por dimensão e requisitos derivados para uma arquitetura de software agêntico}
\label{quadro:1}
\footnotesize
\renewcommand{\arraystretch}{1.25}
\begin{tabularx}{\textwidth}{|X|X|X|}
\hline
\textbf{Dimensão} & \textbf{Lacuna identificada} & \textbf{Requisito derivado} \\
\hline
Indústria e Aplicações & Capacidade técnica sem plano de controle; proliferação de agentes, custo pouco observável e supervisão vagamente definida & R1: coordenar agentes sob um plano de controle único, com custo observável e registro de cada decisão \\
\hline
Desenvolvimento de Software & Avanço concentrado na geração de código; requisitos e projeto pouco atendidos; código plausível, mas estruturalmente frágil & R2: cobrir o ciclo inteiro, da originação da demanda à entrega, e verificar o efeito da mudança, e não apenas o código \\
\hline
Usuário e Experiência & Adaptação com ganho de usabilidade, mas com risco de imprevisibilidade, opacidade e perda de controle & R3: personalizar em áreas delimitadas, de forma explícita, reversível ao padrão e sem gerar código não revisado \\
\hline
Segurança, Ética e Governança & Agentes vulneráveis a manipulação; ação difícil de rastrear e de atribuir & R4: isolar a execução, limitar a produção por contratos verificáveis e reservar a pessoas as decisões que geram consequência \\
\hline
Mercado de Trabalho e Competências & Papéis redefinidos em torno de orquestração e verificação; risco de erosão da competência de quem supervisiona & R5: tornar a supervisão humana informada, com evidências e contexto suficientes para que a decisão seja um julgamento, e não uma formalidade \\
\hline
\end{tabularx}
\fonte{elaborado pelos autores (2026).}
\end{quadro}

Os cinco requisitos convergem para uma escolha de projeto. A via apontada pela evidência não é a do ecossistema totalmente autônomo, nem a da IA reduzida a assistente, mas a da autonomia supervisionada: ampla capacidade de ação dos agentes, subordinada a pontos de decisão humana e a garantias verificáveis por código. É essa escolha que a Arquitetura PASE, apresentada no Capítulo~\ref{cap:4}, procura materializar, depois de o Capítulo~\ref{cap:3} situá-la em relação aos trabalhos que mais se aproximam dela.
